% ============================================================================
% IE 643 Course Project — Intensive Assessment Report
% Team Vision Quest — Early Prediction of Batting Shot-Type from Cricket Videos
%
% >>> READ THIS FIRST <<<
% The instructions award ZERO if the report is AI-generated text or copy-paste
% from papers, and a plagiarism check is run. So:
%   * The structure, tables, figure hooks, math and references below are a
%     SCAFFOLD. They match the required section list exactly.
%   * REWRITE every \TODO{...} and every body paragraph IN YOUR OWN WORDS,
%     describing what YOUR team actually did and found. Do not submit the
%     placeholder prose as-is.
%   * Fill the RESULT numbers from results/temporal_results.json after training.
%
% If the official template (IE643_Project_Report.tex from the Drive folder) has
% a specific \documentclass / packages, paste THESE SECTION BODIES into it.
% This file compiles standalone on Overleaf (pdfLaTeX) as a fallback.
% ============================================================================
\documentclass[11pt]{article}
\usepackage[a4paper,margin=1in]{geometry}
\usepackage{graphicx,booktabs,amsmath,amssymb,hyperref,xcolor,enumitem}
\usepackage{caption}
\usepackage{authblk}
\hypersetup{colorlinks=true,linkcolor=black,citecolor=blue,urlcolor=blue}
\newcommand{\TODO}[1]{\textcolor{red}{[TODO: #1]}}
\newcommand{\NUM}[1]{\textcolor{blue}{#1}}   % fill with real numbers, then remove colour

\title{\textbf{Early Prediction of Batting Shot-Type from Cricket Videos}\\
\large IE 643 Course Project --- Intensive Assessment Report}
\author{Team Vision Quest \\ \TODO{Member 1 Name (Roll No.)} \and \TODO{Member 2 Name (Roll No.)}}
\date{\today}

\begin{document}
\maketitle

% ----------------------------------------------------------------------------
\begin{abstract}
\TODO{Write 150--200 words in your own words.} We address early prediction of a
batter's cricket shot before the stroke is completed. We study two complementary
settings: a video model that predicts the Top-3/Top-5 shot type from only the
first 50\% and 75\% of a clip, and a player-specific pose model that scores shots
as likely/unlikely from the batter's body-keypoint sequence. On the
Cricket Shots -- IPL 2023 dataset (7 classes), our video branch
(R3D-18) reaches a full-clip Top-1 accuracy of \NUM{XX.X\%} and retains
\NUM{XX.X\%} at 50\% observation; the pose branch reaches \NUM{XX.X\%} Top-1.
A temporal analysis shows reliable predictions emerge by \NUM{XX\%} of the clip.
\TODO{State the single most important finding in one sentence.}
\end{abstract}

% ----------------------------------------------------------------------------
\section{Introduction}
\TODO{Rewrite fully.} Anticipating a batter's shot before contact has direct value
for broadcast analytics, coaching feedback and fielding decisions. Human experts
read early cues --- stance, backlift, trigger movement --- well before bat--ball
contact, which motivates asking whether a model can do the same from partial
video. We frame this as \emph{early action prediction}: given only an initial
fraction of a shot clip, predict the final shot label. We also ask whether body
pose alone carries enough early signal, and whether modelling individual batters
helps. \TODO{Explain why deep learning (spatiotemporal CNNs, graph networks on
skeletons) is the right tool here, in your words.}

% ----------------------------------------------------------------------------
\section{Workflow}
\TODO{Convert to your actual timeline.}
\begin{itemize}[nosep]
  \item Problem framing: defined Setting A (video, Top-k) and Setting B (pose, likely/unlikely).
  \item Dataset: adopted Cricket Shots -- IPL 2023 (7 classes); built a stratified train/val/test split.
  \item Preprocessing: uniform 20-frame sampling per clip; MediaPipe BlazePose keypoint extraction.
  \item Modelling: R3D-18 video branch; compact ST-GCN pose branch with optional player embedding.
  \item Early-prediction training: prefix-masking datasets + earliness-aware loss.
  \item Evaluation: Top-k, per-class F1, MRR; temporal analysis across observation ratios.
  \item Interface: Gradio demo for live testing.
\end{itemize}

% ----------------------------------------------------------------------------
\section{Related Work}
\TODO{Write 2--3 paragraphs IN YOUR OWN WORDS summarising these; do not copy abstracts.}
\begin{itemize}[nosep]
  \item \textbf{Early action prediction / anticipation.} \cite{kong2018survey,aliakbarian2017}
        --- observation-ratio protocol and losses that reward early correctness.
  \item \textbf{Spatiotemporal \& skeleton backbones.} \cite{feichtenhofer2019slowfast,yan2018stgcn}
        --- 3D-CNNs for video; ST-GCN for skeleton sequences.
  \item \textbf{Cricket shot classification.} \cite{rajarajeswari2026,sen2021cricshot}
        --- prior work classifies \emph{completed} clips; we extend to the early/pose/player setting.
\end{itemize}

% ----------------------------------------------------------------------------
\section{Methods and Approaches}
\subsection{Work before the prep-presentation review}
\TODO{Describe the literature survey, dataset selection and system design done
pre-prep, in your words.}

\subsection{Work after the prep-presentation review}
\TODO{Describe the implementation and experiments done after prep, in your words.}

\subsection{Setting A --- Early video-based prediction}
Let a clip be sampled into $T$ ordered frames. For observation ratio $g\in(0,1]$
the model sees only the first $k=\lceil gT\rceil$ frames; the remainder are
replaced by a repeat of the $k$-th frame to keep the 3D-CNN input length fixed.
The backbone is R3D-18 (ResNet-18 3D, pretrained on Kinetics-400); its final
layer is replaced by a $C$-class head ($C=7$). \TODO{Explain in your words why a
3D-CNN suits spatiotemporal shot cues.}

\paragraph{Earliness-aware loss.} With cross-entropy $\mathrm{CE}$ and a forced
early view $x_e$ (future frames masked at the smallest training ratio),
\begin{equation}
\mathcal{L} = \mathrm{CE}(f(x), y) + \lambda\, \mathrm{CE}(f(x_e), y),
\end{equation}
so correct predictions from few frames are explicitly rewarded (after
\cite{aliakbarian2017}). We use $\lambda=\NUM{0.5}$, label smoothing $\NUM{0.05}$.

\subsection{Setting B --- Player-specific pose model}
MediaPipe BlazePose gives 33 keypoints per frame, normalised to be
translation/scale invariant (centred on mid-hip, scaled by torso length). A
compact ST-GCN \cite{yan2018stgcn} treats the skeleton as a spatio-temporal
graph; a learnable edge-importance mask weights joint connections. An optional
player-embedding is concatenated before the classifier and is learned only for
batters with $\geq \NUM{15}$ clips, so the model falls back to a global pose
model when player data is scarce. Output logits give Top-k ranking (softmax) and
likely/unlikely decisions (sigmoid $>$ threshold). \TODO{Explain why pose is a
complementary signal to pixels, in your words.}

% ----------------------------------------------------------------------------
\section{Data}
\textbf{Source.} Cricket Shots -- IPL 2023, Kaggle (Apache-2.0),
\url{https://www.kaggle.com/datasets/taarunsridhar/cricket-shots-ipl-2023}.
\textbf{Type.} Short RGB broadcast video clips, one batting shot each.
\textbf{Classes ($C=7$).} \TODO{list exact folder names after \texttt{download --inspect}}:
drive, pull, cut, sweep, flick, slog, misc.
\textbf{Size.} \TODO{total clips = \NUM{N}; per-class counts in Table~\ref{tab:data}.}
\textbf{Split.} Stratified \NUM{70}/\NUM{15}/\NUM{15} train/val/test.
\textbf{Preprocessing.} 20 frames uniformly sampled per clip, resized to
$112\times112$; Kinetics normalisation for the video branch; per-frame BlazePose
keypoints for the pose branch. Partial clips (25/50/75\%) are derived by prefix
masking. \textbf{Player labels} are parsed from filenames where present,
else \texttt{unknown}.

\begin{table}[h]\centering
\caption{Dataset composition (fill from \texttt{download --inspect}).}
\label{tab:data}
\begin{tabular}{lccccccc}
\toprule
Class & drive & pull & cut & sweep & flick & slog & misc \\
\midrule
\# clips & \NUM{--} & \NUM{--} & \NUM{--} & \NUM{--} & \NUM{--} & \NUM{--} & \NUM{--} \\
\bottomrule
\end{tabular}
\end{table}

% ----------------------------------------------------------------------------
\section{Experiments}
\textbf{Code.} \TODO{GitHub link:} \url{https://github.com/<YOUR_USERNAME>/VisionQuest_IE643_CricketShotPrediction}.
\textbf{Hardware.} \TODO{e.g. Google Colab, NVIDIA T4 16\,GB.}
\textbf{Training.} AdamW, lr \NUM{3e-4} (video) / \NUM{9e-4} (pose), weight decay
\NUM{1e-4}, cosine schedule with \NUM{2}-epoch warmup, batch size \NUM{8}
(video), gradient clip \NUM{1.0}, mixed precision, up to \NUM{30} epochs with
early stopping (patience \NUM{7}). The video dataset draws a random observation
ratio from $\{0.5,0.75,1.0\}$ per sample so one model serves all horizons.
\TODO{Report any other settings you actually used.}

% ----------------------------------------------------------------------------
\section{Results}
\TODO{Fill every number from results/temporal\_results.json and insert the saved figures.}

\begin{table}[h]\centering
\caption{Accuracy vs.\ observation ratio (test split).}
\label{tab:acc}
\begin{tabular}{lcccc}
\toprule
Observation ratio & 25\% & 50\% & 75\% & 100\% \\
\midrule
Video Top-1 & \NUM{--} & \NUM{--} & \NUM{--} & \NUM{--} \\
Video Top-3 & \NUM{--} & \NUM{--} & \NUM{--} & \NUM{--} \\
Video Top-5 & \NUM{--} & \NUM{--} & \NUM{--} & \NUM{--} \\
Pose Top-1  & \NUM{--} & \NUM{--} & \NUM{--} & \NUM{--} \\
Pose Top-3  & \NUM{--} & \NUM{--} & \NUM{--} & \NUM{--} \\
\bottomrule
\end{tabular}
\end{table}

\begin{figure}[h]\centering
\includegraphics[width=0.7\linewidth]{../results/figures/accuracy_vs_ratio.png}
\caption{Early-prediction accuracy vs.\ observed clip fraction.}
\label{fig:curve}
\end{figure}

\begin{figure}[h]\centering
\includegraphics[width=0.48\linewidth]{../results/figures/confusion_video_full.png}
\includegraphics[width=0.48\linewidth]{../results/figures/perclass_f1.png}
\caption{Left: confusion matrix (video, full clip). Right: per-class F1.}
\label{fig:cm}
\end{figure}

\textbf{Inferences.} \TODO{In your words: which shots are hardest early and why;
the earliness ratio (\NUM{--}); video vs.\ pose comparison; whether the
player-specific head helped; comparison to the 95.6\% full-clip reference of
\cite{rajarajeswari2026}.}

% ----------------------------------------------------------------------------
\section{Plan for Novelty Assessment}
\TODO{Pick ONE and detail it.} Candidate extensions beyond the brief:
(i) bowling-type conditioning (pace vs.\ spin) as an auxiliary input;
(ii) fusing the video and pose branches (late fusion) and measuring earliness gain;
(iii) an uncertainty-aware early-exit that stops as soon as confidence crosses a
threshold, reporting \emph{how early} per shot. \TODO{State what you will build,
how you will evaluate it, and why it is novel relative to prior work.}

% ----------------------------------------------------------------------------
\section{Conclusion}
\TODO{Rewrite.} We built an end-to-end system for early cricket-shot prediction
in two settings and analysed how accuracy grows with observation. \TODO{State
the main quantitative takeaway and its significance in your own words.}

% ----------------------------------------------------------------------------
\begin{thebibliography}{9}
\bibitem{kong2018survey} Yu Kong and Yun Fu. Human Action Recognition and Prediction: A Survey. \textit{arXiv}, 2018. \url{https://arxiv.org/abs/1806.11230}
\bibitem{aliakbarian2017} Mohammad Sadegh Aliakbarian, Fatemehsadat Saleh, Mathieu Salzmann, Basura Fernando, Lars Petersson and Lars Andersson. Encouraging LSTMs to Anticipate Actions Very Early. \textit{ICCV}, 2017. \url{https://arxiv.org/abs/1703.07023}
\bibitem{feichtenhofer2019slowfast} Christoph Feichtenhofer, Haoqi Fan, Jitendra Malik and Kaiming He. SlowFast Networks for Video Recognition. \textit{ICCV}, 2019. \url{https://arxiv.org/abs/1812.03982}
\bibitem{yan2018stgcn} Sijie Yan, Yuanjun Xiong and Dahua Lin. Spatial Temporal Graph Convolutional Networks for Skeleton-Based Action Recognition. \textit{AAAI}, 2018. \url{https://arxiv.org/abs/1801.07455}
\bibitem{rajarajeswari2026} S. Rajarajeswari, Devika Prashant, Jayanth Srinivasan and Tejaswini Dubey. Deep Learning Based Cricket Batting Shot Classification and Performance Analysis Using Computer Vision. \textit{Scientific Reports}, 2026. \url{https://www.nature.com/articles/s41598-026-52617-1}
\bibitem{sen2021cricshot} Anik Sen, Kaushik Deb, Pranab Kumar Dhar and Takeshi Koshiba. CricShotClassify: An Approach to Classifying Batting Shots from Cricket Videos Using a CNN and GRU. \textit{Sensors}, 2021. \url{https://www.ncbi.nlm.nih.gov/pmc/articles/PMC8072636/}
\bibitem{mediapipe} Valentin Bazarevsky and Ivan Grishchenko. On-Device, Real-Time Body Pose Tracking with MediaPipe BlazePose. \textit{Google Research Blog}, 2020. \url{https://research.google/blog/on-device-real-time-body-pose-tracking-with-mediapipe-blazepose/}. Accessed on: 25th October, 2026.
\bibitem{kaggleds} Taarun Sridhar. Cricket Shots -- IPL 2023 [dataset]. \url{https://www.kaggle.com/datasets/taarunsridhar/cricket-shots-ipl-2023}. Accessed on: 25th October, 2026.
\end{thebibliography}

\end{document}
